\section{Related Work}
\label{sec:related_work}

\subsection{Diffusion-based Lip Sync}
Diff2Lip \cite{mukhopadhyay2024diff2lip} and DrivenVideoEditing \cite{bigioi2024speech} are both end-to-end lip sync methods based on pixel-space audio-conditioned diffusion models. MyTalk \cite{yu2024make} uses diffusion models in the first stage to complete audio-to-motion conversion and uses a VAE \cite{kingma2013auto} in the second stage for motion-to-image generation. StyleSync \cite{zhong2024style} uses transformers in the first stage to convert audio to motion and employs diffusion models in the second stage for motion-to-image generation. DiffDub \cite{liu2024diffdub} uses diffusion autoencoders \cite{preechakul2022diffusion} to convert the masked images into semantic latent codes in the first stage, and uses diffusion models to generate an image conditioned on the semantic latent codes and audio in the second stage.

\subsection{Non-diffusion-based Lip Sync}
Wav2Lip \cite{prajwal2020lip} is the most classic lip sync method that introduced using a pretrained SyncNet \cite{chung2017out} to supervise the training of the lip-sync generator. \cite{gupta2023towards} trains a VQ-VAE \cite{esser2021taming} to encode faces and head poses, and then trains the lip sync generator in the quantized space to generate high-resolution images. StyleSync \cite{guan2023stylesync} follows the overall framework of Wav2Lip, with its main innovation being the use of StyleGAN2 \cite{karras2020analyzing} as the generator backbone. VideoReTalking \cite{cheng2022videoretalking} divides lip sync into three components: semantic-guided reenactment network, lip sync network, and identity-aware refinement and enhancement. DINet \cite{zhang2023dinet} deforms the feature maps to generate mouth shapes conditioned on the driving audio. MuseTalk \cite{zhang2024musetalk} uses the architecture of Stable Diffusion \cite{rombach2022high} for inpainting, but it does not perform diffusion process and uses a discriminator for adversarial learning \cite{goodfellow2014generative}, making it more like a GAN-based framework.

\subsection{Audio-driven Portrait Animation}
Many people may confuse lip sync with audio-driven portrait animation. These two tasks have some similarities but are actually completely different tasks. Lip sync is based on video-to-video editing framework, which needs to keep areas other than the mouth the same as in the input video. Audio-driven portrait animation is based on image-to-video animation framework, which can change the head movement and even facial expressions. Although there are already some audio-driven portrait animation methods based on audio-conditioned LDMs, such as \cite{tian2024emo,chen2024echomimic,xu2024hallo,jiang2024loopy}, they cannot be directly applied to the lip sync task due to the shortcut learning problem.